\documentclass{article}
\usepackage[T2A]{fontenc}
\usepackage[utf8]{inputenc}
\usepackage[russian]{babel}
\usepackage{amsmath,amssymb}
\usepackage{graphicx}
\title{Многоуровневая GRA-обнулёнка для похудения без чувства голода\\ (SlimFit GRA)}
\author{О. Б. Битсоев}
\date{\today}

\begin{document}
\maketitle

\begin{abstract}
Предлагается применение трёхуровневой GRA-обнулёнки (статическая, циклическая, хаотическая) для синтеза персонализированного плана питания и активности, обеспечивающего устойчивое снижение веса при минимальном субъективном голоде. Статическая пена отвечает за достижение целевого веса и сохранение безжировой массы, циклическая — за подавление суточных пиков голода, а хаотическая — за предотвращение метаболических плато и срывов. Оптимизация производится генетическим алгоритмом, использующим упрощённую модель метаболической динамики и предиктор голода. Экспериментальные расчёты показывают принципиальную возможность обнуления всех трёх пен и получения плавной, комфортной траектории похудения.
\end{abstract}

\section{Введение}
\dots
\section{Метод}
\dots
\section{Результаты}
\dots
\section{Заключение}
\dots

\bibliographystyle{plain}
\bibliography{references}
\end{document}
